\hypertarget{functional-programming-modules}{%
\chapter{Functional Programming
Modules}\label{functional-programming-modules}}

Python is not a purely functional language, but it provides powerful
tools for functional programming patterns. These modules are almost
entirely implemented in C, offering near-native performance for
higher-order operations.

\hypertarget{itertools-infinite-and-combinatoric-iterators}{%
\subsection{\texorpdfstring{34.1 \texttt{itertools}: Infinite and
Combinatoric
Iterators}{34.1 itertools: Infinite and Combinatoric Iterators}}\label{itertools-infinite-and-combinatoric-iterators}}

The \passthrough{\lstinline!itertools!} module provides functions that
create iterators for efficient looping. They are ``lazy''they produce
values only when requested, consuming minimal memory.

\hypertarget{infinite-iterators}{%
\subsubsection{1. Infinite Iterators}\label{infinite-iterators}}

\begin{itemize}
\tightlist
\item
  \passthrough{\lstinline!count(start, step)!}: An infinite sequence of
  numbers.
\item
  \passthrough{\lstinline!cycle(iterable)!}: Saves a copy of the
  iterable and repeats it indefinitely.
\item
  \passthrough{\lstinline!repeat(object, times)!}: Returns the same
  object over and over.
\end{itemize}

\hypertarget{combinatoric-iterators}{%
\subsubsection{2. Combinatoric Iterators}\label{combinatoric-iterators}}

\begin{itemize}
\tightlist
\item
  \passthrough{\lstinline!product()!}: Cartesian product (nested
  for-loops).
\item
  \passthrough{\lstinline!permutations()!} /
  \passthrough{\lstinline!combinations()!}: High-performance
  combinatorial generation.
\item
  \textbf{Internals}: These are implemented as C-level classes that
  maintain the current state of the iteration in their internal structs,
  avoiding the overhead of Python-level generators.
\end{itemize}

\hypertarget{functools-higher-order-functions}{%
\subsection{\texorpdfstring{34.2 \texttt{functools}: Higher-Order
Functions}{34.2 functools: Higher-Order Functions}}\label{functools-higher-order-functions}}

The \passthrough{\lstinline!functools!} module is for functions that act
on or return other functions.

\hypertarget{partialfunc-args-keywords}{%
\subsubsection{\texorpdfstring{1.
\texttt{partial(func,\ *args,\ **keywords)}}{1. partial(func, *args, **keywords)}}\label{partialfunc-args-keywords}}

\passthrough{\lstinline!partial!} returns a new
\passthrough{\lstinline!partial!} object (a C struct). *
\textbf{Internals}: It stores the function, the positional arguments,
and the keyword arguments. When called, it merges the stored arguments
with the new ones and calls the original function. This is significantly
faster than using a lambda for the same purpose because it avoids
creating a new Python function object and closure.

\hypertarget{lru_cachemaxsize128-typedfalse}{%
\subsubsection{\texorpdfstring{2.
\texttt{lru\_cache(maxsize=128,\ typed=False)}}{2. lru\_cache(maxsize=128, typed=False)}}\label{lru_cachemaxsize128-typedfalse}}

The Least Recently Used (LRU) cache is implemented using a
\textbf{Dictionary} and a \textbf{Doubly Linked List}. *
\textbf{Dictionary}: Maps the function arguments (hashable) to the
result. * \textbf{Linked List}: Tracks the order of access. When a hit
occurs, the item is moved to the front. If the cache is full, the item
at the tail is evicted. * \textbf{Thread Safety}: It uses a reentrant
lock to ensure that the internal linked list remains consistent across
multiple threads.

\hypertarget{operator-exposing-c-level-opcodes}{%
\subsection{\texorpdfstring{34.3 \texttt{operator}: Exposing C-Level
Opcodes}{34.3 operator: Exposing C-Level Opcodes}}\label{operator-exposing-c-level-opcodes}}

The \passthrough{\lstinline!operator!} module exports a set of efficient
functions corresponding to Python's intrinsic operators. *
\passthrough{\lstinline!operator.add(x, y)!} is equivalent to
\passthrough{\lstinline!x + y!}. *
\passthrough{\lstinline!operator.itemgetter(index)!} is equivalent to
\passthrough{\lstinline!lambda obj: obj[index]!}. *
\passthrough{\lstinline!operator.attrgetter(attr)!} is equivalent to
\passthrough{\lstinline!lambda obj: getattr(obj, attr)!}.

\textbf{Godhood Tip}: Always use
\passthrough{\lstinline!operator.itemgetter!} or
\passthrough{\lstinline!attrgetter!} for sorting or mapping instead of
lambdas. They are implemented in C and avoid the overhead of a Python
function call for every element in the collection.

\hypertarget{singledispatch-generic-functions}{%
\subsection{\texorpdfstring{34.4 \texttt{singledispatch}: Generic
Functions}{34.4 singledispatch: Generic Functions}}\label{singledispatch-generic-functions}}

\passthrough{\lstinline!functools.singledispatch!} allows for function
overloading based on the type of the first argument. *
\textbf{Internals}: It maintains a registry (a dictionary) mapping types
to implementation functions. When the generic function is called, it
performs a lookup in the registry. It also handles inheritance by
traversing the MRO of the input type to find the closest registered
handler.

\begin{center}\rule{0.5\linewidth}{0.5pt}\end{center}

\begin{center}\rule{0.5\linewidth}{0.5pt}\end{center}
